% If your main document uses the article class, replace \chapter with \section.
\chapter{Methodology and Implementation}
\label{chap:methodology_implementation}

This chapter combines the methodological design and the concrete software realization of our benchmark.
The central objective is to compare two quantum toolchains (Bloch and Qiskit) under a matched QAOA workflow:
identical problem instance, identical circuit ansatz, identical classical optimization routine, and identical shot budgets.

\section{Problem Setup}
\label{sec:problem_setup}

We study the MaxCut problem on the 8-node cycle graph \(C_8\), with vertex set
\(V = \{0,\dots,7\}\) and edge set
\[
E = \{(0,1),(1,2),(2,3),(3,4),(4,5),(5,6),(6,7),(7,0)\}.
\]
For a bitstring \(z \in \{0,1\}^{8}\), the cut value is
\[
C(z) = \sum_{(i,j)\in E} (z_i \oplus z_j),
\]
where \(\oplus\) denotes XOR.
The optimization objective is to maximize the expected cut value under a parameterized quantum state.

\section{QAOA Formulation}
\label{sec:qaoa_formulation}

We use depth \(p=1\) QAOA, following the standard construction introduced by Farhi, Goldstone, and Gutmann (2014).
Two real parameters are optimized:
\[
\gamma \in \mathbb{R}, \quad \beta \in \mathbb{R}.
\]

\subsection{Quantum State Preparation and Cost/Mixer Layers}

For \(n=8\) qubits, each evaluation performs:
\begin{enumerate}
    \item Initialize \(\ket{+}^{\otimes n}\) with Hadamard gates on all qubits.
    \item Apply a cost unitary for every edge \((i,j)\in E\), implemented as
    \[
    \mathrm{CX}(i,j)\;\rightarrow\; \mathrm{RZ}_j(-\gamma)\;\rightarrow\; \mathrm{CX}(i,j).
    \]
    \item Apply mixer rotations on all qubits:
    \[
    \mathrm{RX}(2\beta).
    \]
    \item Measure all qubits in the computational basis.
\end{enumerate}

This gate sequence is identical in Bloch and Qiskit implementations.

\subsection{Objective Estimation from Samples}

Given measurement counts at parameters \((\gamma,\beta)\), the empirical probability of bitstring \(z\) is
\[
\hat{p}_{\gamma,\beta}(z)=\frac{\mathrm{counts}(z)}{S},
\]
where \(S\) is the number of shots.
The empirical objective is
\[
\hat{F}(\gamma,\beta)=\sum_{z\in\{0,1\}^{8}} \hat{p}_{\gamma,\beta}(z)\, C(z).
\]
In this project, default shot count is \(S=8192\) for both toolchains.

\section{Classical Optimization Strategy}
\label{sec:classical_optimization}

The classical optimizer is gradient ascent using parameter-shift evaluations.
For one iteration at current \((\gamma,\beta)\):
\begin{align}
\frac{\partial \hat{F}}{\partial \gamma}
&\approx \frac{1}{2}\left[\hat{F}\!\left(\gamma+\frac{\pi}{2},\beta\right)-\hat{F}\!\left(\gamma-\frac{\pi}{2},\beta\right)\right], \\
\frac{\partial \hat{F}}{\partial \beta}
&\approx \hat{F}\!\left(\gamma,\beta+\frac{\pi}{4}\right)-\hat{F}\!\left(\gamma,\beta-\frac{\pi}{4}\right).
\end{align}

Parameters are updated as
\begin{align}
\gamma &\leftarrow \gamma + \eta \frac{\partial \hat{F}}{\partial \gamma}, \\
\beta  &\leftarrow \beta  + \eta \frac{\partial \hat{F}}{\partial \beta},
\end{align}
with learning rate \(\eta=0.12\) in the default configuration.

Each optimization iteration requires five objective evaluations:
\[
\hat{F}(\gamma,\beta),\;
\hat{F}\!\left(\gamma+\frac{\pi}{2},\beta\right),\;
\hat{F}\!\left(\gamma-\frac{\pi}{2},\beta\right),\;
\hat{F}\!\left(\gamma,\beta+\frac{\pi}{4}\right),\;
\hat{F}\!\left(\gamma,\beta-\frac{\pi}{4}\right).
\]
With 25 iterations, this is \(25\times 5=125\) evaluations per toolchain, followed by one final post-optimization evaluation.

\section{Implementation Architecture}
\label{sec:implementation_architecture}

The repository is intentionally split into three executable artifacts:
\begin{enumerate}
    \item \texttt{c8\_maxcut\_qaoa\_p1.bloch}: standalone Bloch circuit with fixed \((\gamma,\beta)\), producing a histogram without optimization.
    \item \texttt{c8\_maxcut\_qaoa\_p1.py}: standalone Python/Qiskit circuit with fixed \((\gamma,\beta)\), producing a histogram without optimization.
    \item \texttt{run\_experiment.py}: complete cross-toolchain experiment runner containing
    all optimization and evaluation logic for both Bloch and Qiskit.
\end{enumerate}

\subsection{Single-Run Scripts}

Both single-run scripts are used as direct sanity checks:
\begin{itemize}
    \item Bloch command: \texttt{bloch c8\_maxcut\_qaoa\_p1.bloch --shots=8192}
    \item Python command: \texttt{poetry run python c8\_maxcut\_qaoa\_p1.py}
\end{itemize}
These scripts avoid optimizer-side effects and expose only the measurement distribution for fixed parameters.

\subsection{Runner Workflow}

\texttt{run\_experiment.py} executes the following steps:
\begin{enumerate}
    \item Configure hyperparameters (defaults: 25 iterations, 8192 shots, \(\eta=0.12\), \(\gamma_0=0.7\), \(\beta_0=0.3\)).
    \item Optimize Bloch parameters with parameter-shift gradient ascent.
    \item Optimize Qiskit parameters with the same routine.
    \item Perform final evaluation for each toolchain at its optimized parameters.
    \item Report matched metrics:
    optimized \((\gamma,\beta)\), optimization runtime, final evaluation runtime, expected cut value, and outcome histograms.
\end{enumerate}

\subsection{Toolchain-Specific Evaluation Details}

\begin{itemize}
    \item \textbf{Bloch path:} the runner updates \((\gamma,\beta)\) literals in the Bloch source, executes the CLI, and parses histogram rows from standard output.
    \item \textbf{Qiskit path:} the runner constructs the same circuit in Python, computes exact statevector probabilities, and samples outcomes using NumPy RNG with explicit seeds.
\end{itemize}

This asymmetry is implementation-level (CLI sampling versus local statevector sampling), but circuit semantics and optimizer protocol are matched.

\section{Reproducibility and Fairness Controls}
\label{sec:reproducibility_fairness}

To preserve comparability:
\begin{itemize}
    \item identical graph instance and edge ordering are used in both toolchains;
    \item identical ansatz structure and gate sequence are enforced;
    \item identical optimizer hyperparameters are used in the shared runner;
    \item identical shot count (\(8192\) by default) is used for both toolchains;
    \item deterministic seed handling is used for Qiskit sampling in optimization and final evaluation.
\end{itemize}

The runner also restores the original Bloch source after execution, so fixed-parameter standalone runs remain stable and repeatable.

\section{Summary}
\label{sec:method_impl_summary}

By unifying methodological assumptions and implementation details in one chapter, the benchmark design becomes explicit and auditable.
The chapter demonstrates that the comparison is not between different algorithms, but between two realizations of the same QAOA procedure under a shared optimization protocol.
